% 问题一：总体路线

本题本质是一个\textbf{区间交叠检测}问题，数据规模较小（$C_{150}^{2} = 11175$ 对），暴力枚举即可在秒级完成。为确保结果正确性，采用\textbf{多种算法交叉验证}的策略。总体路线如下：

\begin{enumerate}
    \item \textbf{数据预处理}：读取附件 1 中 150 个用频计划，解析为结构化数据，预先计算每个计划的所有时间实例列表，避免后续重复计算。
    \item \textbf{多方法检测}：分别采用暴力枚举、频段排序优化、时间排序优化、区间树、扫描线、网格索引、空间哈希等 7 种算法独立检测冲突。
    \item \textbf{一致性校验}：对比 7 种方法的检测结果，确认输出完全一致，确保结果的正确性和可靠性。
    \item \textbf{结果汇总}：记录所有冲突对，统计冲突总数及各类别分布。
\end{enumerate}

考虑到本题数据规模为 $O(10^2)$ 级别，每对计划最多检查 $12 \times 12 = 144$ 次时间交叠判断，总计算量约 $11175 \times 144 \approx 160$ 万次，暴力枚举完全可行。多种方法的交叉验证机制保证了结果的可信度。
